\section*{Erd\H{o}s Problem \#771: a subset avoiding one prescribed sum}
\subsection*{Statement and outcome}
Let \(f(n)\) be the largest integer such that for every \(m\geq1\), some \(S\subseteq[1,n]\) of size \(f(n)\) has no subset summing to \(m\). The asserted asymptotic is
\[
f(n)=\left(\frac12+o(1)\right)\frac n{\log n},
\]
with natural logarithms. We prove the lower bound and the precise arithmetic reduction for the upper bound, taking one explicitly stated Alon--Freiman subset-sum theorem as external input. Attempt status: solved relative to that theorem and the prime number theorem.

\subsection*{A uniform lower bound}
Write \(T=n(n+1)/2\). If \(m>T\), the whole interval \([1,n]\) avoids \(m\). Otherwise let \(p\) be the least prime not dividing \(m\). Such a prime exists. All primes smaller than \(p\) divide \(m\), so
\[
\sum_{\ell<p,\ \ell\ {\rm prime}}\log\ell\leq\log m\leq2\log n+O(1).
\]
The prime number theorem in the form \(\vartheta(x)\sim x\) implies uniformly that
\(p\leq(2+o(1))\log n\). More explicitly, for each \(\varepsilon>0\), the product of primes at most \((2+\varepsilon)\log n\) exceeds \(T\) for large \(n\), so one of them fails to divide \(m\).
Take all multiples of \(p\) up to \(n\). Every subset sum is divisible by \(p\), whereas \(m\) is not. Its size is
\[
\left\lfloor\frac np\right\rfloor\geq
\left(\frac12-o(1)\right)\frac n{\log n}.
\]

\subsection*{The external subset-sum input}
For \(F(n,m)=\max\{|S|:S\subseteq[1,n],\,m\notin\Sigma(S)\}\), let \(s(m)\) be the least positive integer not dividing \(m\). Alon and Freiman's Theorem 1.2 states that, for every fixed \(\varepsilon>0\) and sufficiently large \(n\),
\[
3n^{5/3+\varepsilon}<m<\frac{n^2}{20(\log n)^2}
\quad\Longrightarrow\quad
F(n,m)=\left\lfloor\frac n{s(m)}\right\rfloor+s(m)-2. \tag{1}
\]
The analytic subset-sum theorem behind (1) is not reproved here. We verify all scale and divisibility requirements needed to apply it.

\subsection*{Choosing the difficult sum}
Put \(B=n^2/(40(\log n)^2)\), \(L_j=\operatorname{lcm}(1,\ldots,j)\), and choose the largest \(j\) with \(L_j\leq B\). Set \(m=L_j\). Then \(L_{j+1}>B\), so \(j+1\nmid m\), whereas every integer at most \(j\) divides \(m\). Therefore
\[
s(m)=j+1. \tag{2}
\]
Also \(L_{j+1}/L_j\leq j+1\), whence
\[
\frac B{j+1}<m\leq B. \tag{3}
\]
The identity \(\log L_j=\psi(j)\), and the prime number theorem \(\psi(j)\sim j\), show
\[
j=(2+o(1))\log n.
\]
To check this carefully, maximality gives
\(\psi(j)\leq\log B<\psi(j+1)\); \(j\to\infty\), and both endpoints are asymptotic to \(j\). Thus \(j\sim\log B\sim2\log n\).

Choose, for example, \(\varepsilon=1/12\). By (3),
\(m\gg n^2/(\log n)^3>3n^{7/4}\) for large \(n\), while \(m\leq B<n^2/(20(\log n)^2)\). Hence (1) applies. Since \(f(n)=\min_{m\geq1}F(n,m)\),
\[
f(n)\leq F(n,m)
=\frac n{(2+o(1))\log n}+O(\log n)
=\left(\frac12+o(1)\right)\frac n{\log n}.
\]
Together with the uniform lower bound, this proves the claim.

\subsection*{A direct check of the extremal modular mechanism}
The least nondivisor \(s=s(m)\) is necessarily a prime power: if it had two coprime proper factors larger than one, both would divide \(m\), forcing \(s\mid m\). Write \(m\equiv i\pmod s\), \(1\leq i<s\). Take all multiples of \(s\) up to \(n\), together with \(i-1\) distinct integers congruent to \(1\pmod s\) and \(s-i-1\) congruent to \(-1\pmod s\). For large \(n\) in the above range there are enough choices. For \(s=2\) both extra groups are empty; otherwise they are disjoint. Any subset of the extra integers has residue represented by an integer in
\([-(s-i-1),i-1]\), which contains no integer congruent to \(i\pmod s\). This supplies exactly \(\lfloor n/s\rfloor+s-2\) elements avoiding \(m\), and independently checks the lower side of (1).

\subsection*{Attribution and scope}
The prime-multiple lower construction is Erd\H{o}s--Graham's; the difficult-sum choice and (1) are due to N. Alon and G. Freiman, \emph{On sums of subsets of a set of integers}, Combinatorica 8 (1988), 297--306, Theorem 1.2 and Section 4, \url{https://web.math.princeton.edu/~nalon/PDFS/Publications/On%20sums%20of%20subsets%20of%20a%20set%20of%20integers.pdf}. The exact dependence on that substantial external theorem is part of this completion claim. Statement checked at \url{https://www.erdosproblems.com/771} on 24 September 2026. No novelty or local Lean verification is claimed.
\subsection*{COMPLETION ESTIMATE}
\textbf{COMPLETION: 100\%}
